\documentclass{oxmathproblems}

\printanswers

\course{Understanding Analysis, 2nd Edition}
\sheettitle{Chapter 7 - Section 7.2 - The Riemann Integral} 

\begin{document}
\begin{questions}

%------------------------- Problem 1 -------------------------
\miquestion Let $f$ be a bounded function on $[a, b]$, and let $P$ be an arbitrary partition of $[a, b]$. First, explain why
$U(f) \geq L(f, P)$. Now, prove Lemma 7.2.6: for any bounded function $f$ on $[a, b]$, it is always the case that $U(f) \geq L(f)$.

\begin{solution}
  We'll prove it by contradiction. Suppose that $U(f) < L(f, P)$. By definition of infimum, there must exists a partition $P_{2}$ such that
  $U(f, P_{2}) < L(f, P)$, for otherwise the infimum would've been at least $L(f, P)$.\\
  But $U(f, P_{2}) < L(f, P)$ contradicts Lemma 7.2.4: $L(f, P_{1}) \leq U(f, P_{2})$ for any partitions $P_{1}, P_{2}$.

  To use this fact to prove Lemma 7.2.6, note that $L(f)$ is the supremum of all lower sums while $U(f)$ is an upper bound. Since by definition,
  an upper bound is \textit{at least} the supremum, it follows that $U(f) \geq L(f)$.

  Alternatively, we can prove Lemma 7.2.6 directly by contradiction as well. Suppose $U(f) < L(f)$ and let $m$ be the midpoint i.e. $m=\frac{L(f)+U(f)}{2}$.
  By definition of infimum, there exists a partition $P_{1}$ such that $U(f, P_{1}) < m$. By definition of supremum, there exists a
  partition $P_{2}$ such that $m < L(f, P_{2})$. But $U(f, P_{1}) < m < L(f, P_{2})$ contradicts Lemma 7.2.4.
\end{solution}

\end{questions}
\end{document}
